% Roboczy podział AF 23: etykiety historyczne pozostają niezmienione.
\setcounter{section}{36}
\section{Twierdzenie o jądrze Schwartza}
\label{sec:dt-schwartz-kernel}
\providecommand{\atlasref}[2]{#2}

Korzystamy z własności funkcji testowych, przestrzeni
dystrybucji i ciągłości form na etapach granicy induktywnej.
Najpierw udowodnimy rozkład funkcji testowych na szeregi
tensorów elementarnych, a następnie twierdzenie o jądrze
w wersji operatorowej. Dowód nie zakłada jeszcze
identyfikacji ukończonego iloczynu induktywnego funkcji
testowych z przestrzenią testową na iloczynie obszarów.

\subsection{Rozkład tensorowy i reprezentacja operatorów}

\noindent\emph{Powiązania tematyczne w Bibliotece matematycznej:} \atlasref{LA-10::thm:lin10-istnienie}{LA-10: algebraiczna konstrukcja iloczynu tensorowego (przed uzupełnieniem topologicznym)}.


\begin{definicja}[Forma dwuliniowa zadana jądrem]
\label{def:kernel-pair23}
Niech $\Omega\subseteq\R^{d_1}$ i
$\Lambda\subseteq\R^{d_2}$ będą niepustymi zbiorami otwartymi.
Dla $K\in\mathcal D'(\Omega\times\Lambda)$ definiujemy
formę na parach funkcji testowych przez
$b_K(f,g)=\langle K,f\otimes g\rangle$,
dla $f\in\mathcal D(\Omega)$ i $g\in\mathcal D(\Lambda)$,
gdzie $(f\otimes g)(x,y)=f(x)g(y)$.
Parowanie jest dwuliniowe także nad $\C$; nie stosujemy
sprzężenia zespolonego.
\end{definicja}

\begin{lemat}[Rozkład testu na szeregi tensorowe]
\label{lem:tensor-density23}
Niech $d_1,d_2\ge1$, a $\Omega\subseteq\R^{d_1}$ i
$\Lambda\subseteq\R^{d_2}$ będą niepustymi zbiorami otwartymi.
Dla każdego $\Phi\in\mathcal D(\Omega\times\Lambda)$
istnieje przedstawienie $\Phi$ jako sumy skończenie
wielu szeregów postaci
$\sum_{n\in\mathbb Z^{d_2}}f_n\otimes g_n$,
zbieżnych w $\mathcal D(\Omega\times\Lambda)$,
gdzie wszystkie $f_n$ mają nośniki w jednym zwartym
$K\Subset\Omega$, wszystkie $g_n$ w jednym zwartym
$L\Subset\Lambda$ i dla każdego $m,r\in\N_0$
\[
 \sum_n p_{K,m}(f_n)\,p_{L,r}(g_n)<\infty.
\]
W szczególności algebraiczna przestrzeń tensorowa
$\mathcal D(\Omega)\otimes\mathcal D(\Lambda)$
jest gęsta w $\mathcal D(\Omega\times\Lambda)$.
\end{lemat}
\begin{proof}
Wybieramy lokalnie skończone gładkie rozkłady jedności
$(\eta_i)$ na $\Omega$ i $(\theta_j)$ na $\Lambda$
podporządkowane pokryciom stosunkowo zwartymi otwartymi
sześcianami. Dla każdego $i,j$ ustalamy zbiory zwarte
$K_i=\supp\eta_i$ i $L_j=\supp\theta_j$ oraz otwarty
sześcian $V_j$ z $L_j\subset V_j\Subset\Lambda$.
Wybieramy rzeczywistą $\chi_j\in\mathcal D(V_j)$,
równą $1$ w otoczeniu $L_j$ (Lemat~\ref{lem:cutoffs17}).
Funkcję $\Phi_{ij}=\eta_i\theta_j\Phi$ przedłużamy
zerem w zmiennej $y$ poza $V_j$ i okresowo z jego
domknięcia $Q_j$; przedłużenie jest gładkie,
gdyż nośnik w $y$ leży w $L_j$.
Niech $a_j$ będzie środkiem $Q_j$, a $2R_j$ jego
długością boku. Kładziemy
\[
 e_{jn}(y)=e^{i\pi n\cdot(y-a_j)/R_j},
 \quad
 f_{ijn}(x)=\frac{1}{(2R_j)^{d_2}}
 \int_{Q_j}\Phi_{ij}(x,y)\overline{e_{jn}(y)}
 \,\mathrm d\lambda_{d_2}(y),
 \quad g_{jn}=\chi_j e_{jn}.
\]
Funkcje $f_{ijn}$ są gładkie i mają nośniki w $K_i$:
pochodne względem $x$ obliczamy pod całką po zwartym
sześcianie. Niech $L_j^*=\supp\chi_j\cup L_j$.
Dla $m,r,t\in\N_0$, całkując przez części za pomocą
$(1-\Delta_y)^t$, otrzymujemy
\[
 \begin{aligned}
 p_{K_i,m}(f_{ijn})&\le C_{j,t}
 (1+|n|)^{-2t}p_{K_i\times L_j^*,m+2t}(\Phi_{ij}),\\
 p_{L_j^*,r}(g_{jn})&\le C_{j,r}'(1+|n|)^r .
 \end{aligned}
\]
Pierwsza nierówność wynika z tożsamości
$(1-\Delta_y)e_{jn}=(1+\pi^2|n|^2/R_j^2)e_{jn}$,
braku wyrazów brzegowych oraz ograniczenia pochodnych
funkcji podcałkowej; druga ze wzoru Leibniza.
Przy $2t>r+d_2$ szereg iloczynów seminorm jest zbieżny
(liczba punktów kratowych o normie maksimum $s$
jest rzędu co najwyżej $(1+s)^{d_2-1}$).
W szczególności dla każdego $q$ szereg
$\sum_n f_{ijn}\otimes g_{jn}$ jest bezwzględnie
zbieżny w seminormie $p_{K_i\times L_j^*,q}$,
stosując powyższe oszacowanie z $m=r=q$.
Zupełność etapu $\mathcal D_{K_i\times L_j^*}$
(Twierdzenie~\ref{thm:DK-frechet18}) daje granicę.

Dla ustalonego $x$ szereg $\sum_n f_{ijn}(x)e_{jn}(y)$
na $Q_j$ jest jednostajnie zbieżny, a jego granicę
identyfikujemy z $\Phi_{ij}(x,\cdot)$ dzięki
zupełności układu trygonometrycznego w $L^2(Q_j)$
i ciągłości obu funkcji. Po pomnożeniu przez $\chi_j$
równość obowiązuje na całej $\Lambda$: na $Q_j$
mamy $\chi_j\Phi_{ij}=\Phi_{ij}$, a poza $Q_j$
obie strony znikają. Warunek $\supp\chi_j\subset V_j$
jest istotny: eliminuje inne kopie okresowe.
Stąd $\Phi_{ij}=\sum_n f_{ijn}\otimes g_{jn}$.
Tylko skończenie wiele par $(i,j)$ daje niezerowe
$\Phi_{ij}$, dzięki lokalnej skończoności i zwartości
projekcji nośnika $\Phi$. Sumujemy te rozwinięcia,
biorąc sumy mnogościowe odpowiednich zwartych zbiorów
$K_i$ i $L_j^*$. Daje to wspólny zwarty nośnik
i żądaną gęstość.

Nad $\R$ w powyższych szeregach dla rzeczywistej $\Phi$
zastępujemy każdy składnik przez jego część rzeczywistą:
$\operatorname{Re}f_{ijn}\otimes\operatorname{Re}g_{jn}
-\operatorname{Im}f_{ijn}\otimes\operatorname{Im}g_{jn}$.
Suma pozostaje zbieżna z tymi samymi oszacowaniami
(z dodatkowym czynnikiem $2$).
\end{proof}

\begin{lemat}[Ciągłość formy na ustalonych etapach]
\label{lem:separate-joint23}
Niech $E,F$ będą przestrzeniami Frécheta nad
$\K\in\{\R,\C\}$, a $b:E\times F\to\K$
formą dwuliniową ciągłą osobno w każdej zmiennej.
Wtedy $b$ jest łącznie ciągła.
\end{lemat}
\begin{proof}
Wybierzmy malejącą przeliczalną bazę
$(V_n)$ otoczeń zera w $F$, złożoną ze zbiorów
zrównoważonych. Niech
$A_n=\{x\in E: |b(x,y)|\le1\ \text{dla wszystkich }y\in V_n\}$.
Każdy $A_n$ jest domknięty, ponieważ $x\mapsto b(x,y)$
jest ciągły dla ustalonego $y$. Z ciągłości w drugiej
zmiennej wynika $E=\bigcup_n A_n$.
Z twierdzenia Baire'a pewien $A_n$ ma niepuste wnętrze;
ponieważ $A_n$ jest wypukły, zrównoważony i domknięty,
ma wówczas otoczenie zera we wnętrzu. Istotnie,
jeśli $x_0+U\subseteq A_n$ dla otwartego symetrycznego
otoczenia zera $U$, to również $-x_0+U\subseteq A_n$;
wypukłość daje $U\subseteq A_n$.
Stąd $|b(x,y)|\le1$ na $U\times V_n$.
Skalowanie daje ciągłość w $(0,0)$, a dla
dowolnego $(x_0,y_0)$ korzystamy z rozwinięcia
$b(x_0+x,y_0+y)-b(x_0,y_0)
=b(x_0,y)+b(x,y_0)+b(x,y)$ i ciągłości osobnej.
\end{proof}

\begin{uwaga}[Ciągłość osobna i łączna]
\label{rem:kernel-separate23}
Dla formy $b:\mathcal D(\Omega)\times\mathcal D(\Lambda)\to\K$
ciągłość osobna nie musi implikować łącznej ciągłości na
całym iloczynie przestrzeni LF. Implikacja ta jest jednak
prawdziwa po ograniczeniu formy do
$\mathcal D_K(\Omega)\times\mathcal D_L(\Lambda)$,
gdyż obie przestrzenie są przestrzeniami Frécheta.
To rozróżnienie jest istotne w twierdzeniu o jądrze;
właściwą globalną topologią tensorową jest topologia
induktywna, omówiona w następnym rozdziale.
\end{uwaga}

\begin{twierdzenie}[Twierdzenie o jądrze Schwartza]
\label{thm:schwartz-kernel23}
Niech $\Omega\subseteq\R^{d_1}$ i
$\Lambda\subseteq\R^{d_2}$ będą niepustymi zbiorami
otwartymi. Dla każdego ciągłego operatora liniowego
\[
 A:\mathcal D(\Omega)\longrightarrow
 \mathcal D'(\Lambda)_\beta
\]
istnieje dokładnie jedna dystrybucja
$K_A\in\mathcal D'(\Omega\times\Lambda)$ taka, że
\[
 \langle Af,g\rangle=\langle K_A,f\otimes g\rangle
 \qquad(f\in\mathcal D(\Omega),\ g\in\mathcal D(\Lambda)).
\]
Odwrotnie, każda dystrybucja
$K\in\mathcal D'(\Omega\times\Lambda)$ określa w ten
sposób ciągły operator liniowy
$A_K:\mathcal D(\Omega)\to\mathcal D'(\Lambda)_\beta$.
Indeks $\beta$ oznacza silną topologię przestrzeni dualnej,
czyli jednostajną zbieżność na ograniczonych zbiorach
funkcji testowych.
\end{twierdzenie}
\begin{proof}
Niech $b(f,g)=\langle Af,g\rangle$. Ciągłość $A$
i ciągłość ewaluacji na pojedynczej funkcji testowej
zapewniają osobną ciągłość $b$. Na przestrzeniach
$\mathcal D_{K_x}(\Omega)$ i $\mathcal D_{K_y}(\Lambda)$
forma $b$ jest wspólnie ciągła na mocy
Lematu~\ref{lem:separate-joint23}; zatem istnieją
$C,m,r$ takie, że
$|b(f,g)|\le C p_{K_x,m}(f)p_{K_y,r}(g)$.
Zbudujemy jeden liniowy funkcjonał na całej przestrzeni
testowej, a następnie sprawdzimy jego niezależność
od wyborów pomocniczych. Wybierzmy lokalnie skończone
rozkłady jedności $(\eta_i)$ na $\Omega$ oraz
$(\theta_j)$ na $\Lambda$, których elementy mają zwarte
nośniki w sześcianach stosunkowo zwartych. Produkty
$\eta_i(x)\theta_j(y)$ tworzą rozkład jedności na
$\Omega\times\Lambda$. Dla danej
$\Phi\in\mathcal D(\Omega\times\Lambda)$ tylko skończenie
wiele funkcji $\Phi_{ij}=\eta_i\theta_j\Phi$ jest niezerowych.
Dla każdej pary $(i,j)$ wybieramy z góry odcięcie
$\chi_j$ i układ Fouriera z dowodu
Lematu~\ref{lem:tensor-density23}, niezależnie od $\Phi$.
Otrzymujemy $\Phi_{ij}=\sum_n f_{ijn}\otimes g_{jn}$,
gdzie $f_{ijn}$ zależy liniowo od $\Phi$,
a $g_{jn}$ nie zależy od $\Phi$.
Nad $\R$ stosujemy opisane w lemacie dwa rzeczywiste
składniki; dalej zapisujemy wzory w przypadku zespolonym.
Definiujemy
\[
 \langle K_A,\Phi\rangle=\sum_{i,j}\sum_n b(f_{ijn},g_{jn}),
\]
gdzie suma po $i,j$ jest skończona, a wszystkie szeregi
wewnętrzne są bezwzględnie zbieżne na mocy oszacowań
seminormowych i ciągłości $b$ na odpowiednich etapach.
Konstrukcja jest liniowa, ponieważ wybory odcięć i baz
Fouriera nie zależą od $\Phi$, a współczynniki Fouriera
zależą od niej liniowo.

Aby sprawdzić ciągłość, ustalmy $M\Subset\Omega\times\Lambda$.
Lokalna skończoność rozkładów jedności daje skończoną
liczbę par $(i,j)$, które mogą wystąpić dla
$\Phi\in\mathcal D_M$. Dla każdej takiej pary istnieją $C_{ij},m_{ij},r_{ij}$
z oszacowania dwuliniowego na $\mathcal D_{K_i}
\times\mathcal D_{L_j^*}$. Wybieramy $t_{ij}$ z
$2t_{ij}>r_{ij}+d_2$. Z oszacowań w dowodzie
Lematu~\ref{lem:tensor-density23} i wzoru Leibniza
dla $\eta_i\theta_j$ otrzymujemy
\[
 |b(f_{ijn},g_{jn})|\le C'_{ij}
 (1+|n|)^{r_{ij}-2t_{ij}}
 p_{M,m_{ij}+2t_{ij}}(\Phi).
\]
Pochodne $\Phi$ znikają poza $M$, zatem po prawej
wystarcza seminorma na $M$. Sumujemy po $n$,
a następnie po skończonej liczbie par i bierzemy
$r_M=\max_{i,j}(m_{ij}+2t_{ij})$ (dla pustej
rodziny par funkcjonał jest zerowy). Otrzymujemy
stałą $C_M>0$ i $r_M\in\N_0$ takie, że
\[
 |\langle K_A,\Phi\rangle|\le C_M p_{M,r_M}(\Phi)
 \qquad (\Phi\in\mathcal D_M(\Omega\times\Lambda)).
\]
Jest to ciągłość na każdym etapie, więc własność
uniwersalna granicy induktywnej daje
$K_A\in\mathcal D'(\Omega\times\Lambda)$.

Dla $\Phi=f\otimes g$ współczynniki mają szczególną
postać $f_{ijn}=(\eta_i f)c_{jn}(\theta_j g)$,
gdzie $c_{jn}$ jest skalarnym współczynnikiem Fouriera
w $Q_j$. Z dowodu lematu
$\sum_n c_{jn}(\theta_j g)g_{jn}=\theta_j g$
w ustalonym etapie $\mathcal D_{L_j^*}$.
Liniowość i ciągłość $b$ względem drugiego argumentu
przy ustalonym $\eta_i f$ dają
\[
 \sum_n b(f_{ijn},g_{jn})=b(\eta_i f,\theta_j g).
\]
Nad $\R$ tę samą równość otrzymujemy, rozwijając
rzeczywistą część szeregu i korzystając z rzeczywistej
dwuliniowości. Sumując skończenie wiele składników
rozkładów jedności, dostajemy
$\langle K_A,f\otimes g\rangle=b(f,g)$.
Jeżeli zmienimy rozkłady jedności i odcięcia, otrzymamy
inny ciągły funkcjonał na całej przestrzeni $\mathcal D$,
ale o tych samych wartościach na tensorach elementarnych.
Ich liniowa powłoka jest gęsta na mocy
Lematu~\ref{lem:tensor-density23}; dwa funkcjonały muszą
więc być równe. To zarazem uzasadnia jednoznaczność,
bez zakładania gęstości tensorów o nośnikach w arbitralnym
zwartym zbiorze $M$.

W odwrotną stronę, dla $K\in\mathcal D'(\Omega\times\Lambda)$
forma $b_K(f,g)=\langle K,f\otimes g\rangle$ jest
łącznie ciągła na każdym
$\mathcal D_{K_x}(\Omega)\times\mathcal D_{K_y}(\Lambda)$,
ponieważ dla ustalonych niepustych zwartych $K_x,K_y$
z ciągłości $K$ na etapie $\mathcal D_{K_x\times K_y}$
otrzymujemy pewne $C>0$ i $q\in\N_0$ z
\[
 |\langle K,f\otimes g\rangle|
 \le C p_{K_x\times K_y,q}(f\otimes g)
 \le C p_{K_x,q}(f)p_{K_y,q}(g).
\]
Ostatnia nierówność wynika z rozdzielenia pochodnych
względem $x$ i $y$. Dla ustalonego $f$ wybieramy
$K_x$ zawierający jego nośnik. Forma $g\mapsto b_K(f,g)$
jest liniowa i ciągła na każdym etapie $\mathcal D_{K_y}$,
więc na mocy Twierdzenia~\ref{thm:LF-universal18}
wyznacza $A_Kf\in\mathcal D'(\Lambda)$. Dwuliniowość
zapewnia liniowość $A_K$.
Aby sprawdzić ciągłość w topologii silnej, niech
$B\subseteq\mathcal D(\Lambda)$ będzie ograniczony.
Z Twierdzenia~\ref{thm:D-bounded18} istnieje
$K_y\Subset\Lambda$, dla którego $B$ jest ograniczony
w $\mathcal D_{K_y}$. Na każdym etapie $\mathcal D_{K_x}$
mamy, dzięki wspólnej ciągłości formy,
\[
 \sup_{g\in B}|\langle A_Kf,g\rangle|
 \le C\,p_{K_x,m}(f)\sup_{g\in B}p_{K_y,r}(g).
\]
To dowodzi ciągłości $A_K$ na każdym etapie, a więc,
przez własność uniwersalną granicy induktywnej,
na całej $\mathcal D(\Omega)$. Jednoznaczność $K_A$
wynika ponownie z gęstości tensorów elementarnych.
\end{proof}

\begin{przyklad}[Jądro operatora o gładkim jądrze]
\label{ex:smooth-kernel23}
Jeżeli $k\in C^\infty(\Omega\times\Lambda)$, to dystrybucja
regularna $T_k$ jest jądrem operatora
$(Af)(y)=\int_\Omega k(x,y)f(x)\,\mathrm d\lambda_{d_1}(x)$.
Dla $g\in\mathcal D(\Lambda)$ mamy
\[
\langle Af,g\rangle
=\int_{\Lambda}\int_{\Omega}k(x,y)f(x)g(y)
\,\mathrm d\lambda_{d_1}(x)\,\mathrm d\lambda_{d_2}(y).
\]
Dla ustalonego $f$ różniczkowanie względem $y$
pod całką jest uzasadnione przez ograniczoność pochodnych
$k$ na iloczynie $\supp f$ i dowolnego zwartego
otoczenia punktu $y$; funkcja $Af$ jest gładka.
Twierdzenie Fubiniego stosujemy, ponieważ funkcja
podcałkowa jest ograniczona i znika poza
$\supp f\times\supp g$. Nie wymagamy, żeby $k$
miało zwarty nośnik. Identyfikując $Af$ z dystrybucją
regularną, z twierdzenia o jądrze otrzymujemy
ciągłość operatora $\mathcal D(\Omega)\to
\mathcal D'(\Lambda)_\beta$.
\end{przyklad}

\begin{przyklad}[Jądro osobliwe na przekątnej]
\label{ex:diagonal-kernel23}
Dla $\Omega=\Lambda\subseteq\R^d$ operator naturalnego
włożenia $\mathcal D(\Omega)\hookrightarrow\mathcal D'(\Omega)$
ma jądro $\delta_\Delta$ zadane przez
$\langle\delta_\Delta,\Phi\rangle
=\int_\Omega\Phi(x,x)\,\mathrm d\lambda_d(x)$.
Dla $\Phi\in\mathcal D_M(\Omega\times\Omega)$
całka jest ograniczona przez
$\lambda_d(\operatorname{pr}_1 M)p_{M,0}(\Phi)$,
co dowodzi ciągłości na każdym etapie.
Na tensorze $f\otimes g$ otrzymujemy $\int fg\,
\mathrm d\lambda_d$, czyli parowanie naturalnego
włożenia. Nośnik $\delta_\Delta$ jest zawarty
w przekątnej $\Delta=\{(x,x):x\in\Omega\}$;
każde otoczenie punktu $(x,x)$ zawiera iloczyn
małych kul o środku $x$, na którym można wybrać
nieujemny tensor $f\otimes f$ z $f\ne0$.
Wtedy $\langle\delta_\Delta,f\otimes f\rangle>0$,
więc nośnik jest równy $\Delta$.
Gdyby $\delta_\Delta=T_h$ dla
$h\in L^1_{\mathrm{loc}}(\Omega\times\Omega)$,
to ograniczenie do otwartego dopełnienia $\Delta$
byłoby zerowe. Różnowartościowość przyporządkowania
funkcji dystrybucji regularnej daje $h=0$ prawie
wszędzie poza $\Delta$. Przekątna ma miarę
$\lambda_{2d}$ zero, zatem $h=0$ prawie wszędzie,
co przeczy $\delta_\Delta\ne0$.
\end{przyklad}

\begin{uwaga}[Zakres twierdzenia o jądrze]
\label{rem:kernel-scope23}
Twierdzenie~\ref{thm:schwartz-kernel23} opisuje wszystkie
ciągłe operatory $\mathcal D(\Omega)\to\mathcal D'(\Lambda)_\beta$.
Nie oznacza ono, że jądro jest funkcją ani że operacja
składania dowolnych takich operatorów jest automatycznie
zdefiniowana. Twierdzenia trwałości nuklearności zostały udowodnione
w rozdziale~\ref{sec:montel-nuklearne23}; następny rozdział omawia
induktywny iloczyn tensorowy funkcji testowych. Pełne rozwiązywanie równań
różniczkowych pozostaje w osobnym skrypcie.
\end{uwaga}
